% Options for packages loaded elsewhere
\PassOptionsToPackage{unicode}{hyperref}
\PassOptionsToPackage{hyphens}{url}
\PassOptionsToPackage{dvipsnames,svgnames,x11names}{xcolor}
%
\documentclass[
  11pt,
  a4paper,
]{article}
\usepackage{amsmath,amssymb}
\usepackage{iftex}
\ifPDFTeX
  \usepackage[T1]{fontenc}
  \usepackage[utf8]{inputenc}
  \usepackage{textcomp} % provide euro and other symbols
\else % if luatex or xetex
  \usepackage{unicode-math} % this also loads fontspec
  \defaultfontfeatures{Scale=MatchLowercase}
  \defaultfontfeatures[\rmfamily]{Ligatures=TeX,Scale=1}
\fi
\usepackage{lmodern}
\ifPDFTeX\else
  % xetex/luatex font selection
\fi
% Use upquote if available, for straight quotes in verbatim environments
\IfFileExists{upquote.sty}{\usepackage{upquote}}{}
\IfFileExists{microtype.sty}{% use microtype if available
  \usepackage[]{microtype}
  \UseMicrotypeSet[protrusion]{basicmath} % disable protrusion for tt fonts
}{}
\makeatletter
\@ifundefined{KOMAClassName}{% if non-KOMA class
  \IfFileExists{parskip.sty}{%
    \usepackage{parskip}
  }{% else
    \setlength{\parindent}{0pt}
    \setlength{\parskip}{6pt plus 2pt minus 1pt}}
}{% if KOMA class
  \KOMAoptions{parskip=half}}
\makeatother
\usepackage{xcolor}
\usepackage[margin=2.25cm]{geometry}
\usepackage{color}
\usepackage{fancyvrb}
\newcommand{\VerbBar}{|}
\newcommand{\VERB}{\Verb[commandchars=\\\{\}]}
\DefineVerbatimEnvironment{Highlighting}{Verbatim}{commandchars=\\\{\}}
% Add ',fontsize=\small' for more characters per line
\newenvironment{Shaded}{}{}
\newcommand{\AlertTok}[1]{\textcolor[rgb]{1.00,0.00,0.00}{\textbf{#1}}}
\newcommand{\AnnotationTok}[1]{\textcolor[rgb]{0.38,0.63,0.69}{\textbf{\textit{#1}}}}
\newcommand{\AttributeTok}[1]{\textcolor[rgb]{0.49,0.56,0.16}{#1}}
\newcommand{\BaseNTok}[1]{\textcolor[rgb]{0.25,0.63,0.44}{#1}}
\newcommand{\BuiltInTok}[1]{\textcolor[rgb]{0.00,0.50,0.00}{#1}}
\newcommand{\CharTok}[1]{\textcolor[rgb]{0.25,0.44,0.63}{#1}}
\newcommand{\CommentTok}[1]{\textcolor[rgb]{0.38,0.63,0.69}{\textit{#1}}}
\newcommand{\CommentVarTok}[1]{\textcolor[rgb]{0.38,0.63,0.69}{\textbf{\textit{#1}}}}
\newcommand{\ConstantTok}[1]{\textcolor[rgb]{0.53,0.00,0.00}{#1}}
\newcommand{\ControlFlowTok}[1]{\textcolor[rgb]{0.00,0.44,0.13}{\textbf{#1}}}
\newcommand{\DataTypeTok}[1]{\textcolor[rgb]{0.56,0.13,0.00}{#1}}
\newcommand{\DecValTok}[1]{\textcolor[rgb]{0.25,0.63,0.44}{#1}}
\newcommand{\DocumentationTok}[1]{\textcolor[rgb]{0.73,0.13,0.13}{\textit{#1}}}
\newcommand{\ErrorTok}[1]{\textcolor[rgb]{1.00,0.00,0.00}{\textbf{#1}}}
\newcommand{\ExtensionTok}[1]{#1}
\newcommand{\FloatTok}[1]{\textcolor[rgb]{0.25,0.63,0.44}{#1}}
\newcommand{\FunctionTok}[1]{\textcolor[rgb]{0.02,0.16,0.49}{#1}}
\newcommand{\ImportTok}[1]{\textcolor[rgb]{0.00,0.50,0.00}{\textbf{#1}}}
\newcommand{\InformationTok}[1]{\textcolor[rgb]{0.38,0.63,0.69}{\textbf{\textit{#1}}}}
\newcommand{\KeywordTok}[1]{\textcolor[rgb]{0.00,0.44,0.13}{\textbf{#1}}}
\newcommand{\NormalTok}[1]{#1}
\newcommand{\OperatorTok}[1]{\textcolor[rgb]{0.40,0.40,0.40}{#1}}
\newcommand{\OtherTok}[1]{\textcolor[rgb]{0.00,0.44,0.13}{#1}}
\newcommand{\PreprocessorTok}[1]{\textcolor[rgb]{0.74,0.48,0.00}{#1}}
\newcommand{\RegionMarkerTok}[1]{#1}
\newcommand{\SpecialCharTok}[1]{\textcolor[rgb]{0.25,0.44,0.63}{#1}}
\newcommand{\SpecialStringTok}[1]{\textcolor[rgb]{0.73,0.40,0.53}{#1}}
\newcommand{\StringTok}[1]{\textcolor[rgb]{0.25,0.44,0.63}{#1}}
\newcommand{\VariableTok}[1]{\textcolor[rgb]{0.10,0.09,0.49}{#1}}
\newcommand{\VerbatimStringTok}[1]{\textcolor[rgb]{0.25,0.44,0.63}{#1}}
\newcommand{\WarningTok}[1]{\textcolor[rgb]{0.38,0.63,0.69}{\textbf{\textit{#1}}}}
\setlength{\emergencystretch}{3em} % prevent overfull lines
\providecommand{\tightlist}{%
  \setlength{\itemsep}{0pt}\setlength{\parskip}{0pt}}
\setcounter{secnumdepth}{-\maxdimen} % remove section numbering
\ifLuaTeX
  \usepackage{selnolig}  % disable illegal ligatures
\fi
\usepackage{bookmark}
\IfFileExists{xurl.sty}{\usepackage{xurl}}{} % add URL line breaks if available
\urlstyle{same}
\hypersetup{
  colorlinks=true,
  linkcolor={blue},
  filecolor={Maroon},
  citecolor={Blue},
  urlcolor={blue},
  pdfcreator={LaTeX via pandoc}}

\author{}
\date{}

\usepackage{microtype}
\usepackage{fancyhdr}
\pagestyle{fancy}
\fancyhf{}
\lhead{Documentation-Driven Threat Analysis}
\rhead{BA4 closure}
\cfoot{\thepage}
\setlength{\headheight}{14pt}
\setlength{\emergencystretch}{4em}
\hbadness=10000
\hfuzz=20pt
\sloppy
\begin{document}

\section{DDTA BA4-T3 projection closure review -
R1}\label{ddta-ba4-t3-projection-closure-review---r1}

\textbf{Status:} CLOSED / PASS

\textbf{Repository baseline reviewed:}
\texttt{dcb4605448de4ac5331f10ff090a9f2ab677427e}

\textbf{Executed scope:} BA4-T3 only - integrated projection boundary,
interpretation/coverage and cross-projection closure review.

\subsection{1. Closure question}\label{closure-question}

Can the combined BA4-T1/T2 candidate be reduced further while still
preserving authority, semantic rendering, coverage/completeness,
qualification, method-rule accountability, cross-projection comparison
and change rebuild?

The review attempts to falsify the candidate by deleting or merging
responsibilities before accepting closure.

\subsection{2. Inputs}\label{inputs}

The review uses:

\begin{itemize}
\tightlist
\item
  \texttt{BA4\_PROJECTION\_BOUNDARY\_TRACEABILITY\_CANDIDATE\_R1.md};
\item
  \texttt{BA4\_T1\_PROJECTION\_BOUNDARY\_TRACEABILITY\_SEMANTIC\_PRESERVATION\_R1.md};
\item
  \texttt{BA4\_PROJECTION\_BOUNDARY\_INTERPRETATION\_COVERAGE\_CANDIDATE\_R2.md};
\item
  \texttt{BA4\_T2\_METHOD\_INTERPRETATION\_COVERAGE\_CROSS\_PROJECTION\_R1.md};
\item
  \texttt{DDTA\_CURRENT\_RESEARCH\_STATE\_R15.md};
\item
  closed BA1, BA2 and BA3 contracts; and
\item
  the facial M1-M4, order/WMS and provider-normalization controls
  already used by BA4-T1/T2.
\end{itemize}

No complete threat-method schema, AnalysisRecord/Common Finding or
ThreatForge implementation model is introduced.

\subsection{3. Attack A - remove
coverageMode}\label{attack-a---remove-coveragemode}

Candidate simplification:

\begin{Shaded}
\begin{Highlighting}[]
\NormalTok{selectionCoverageContract:}
\NormalTok{  "current transfer propositions"}
\end{Highlighting}
\end{Shaded}

with no explicit statement of selective vs exhaustive intent.

Two valid consumers can read that same scope as:

\begin{Shaded}
\begin{Highlighting}[]
\NormalTok{A: show every current transfer}
\NormalTok{B: show only transfers useful to one analysis question}
\end{Highlighting}
\end{Shaded}

The omission of one eligible transfer therefore cannot be classified as
either a defect or an intentional omission.

\textbf{Disposition:} removal \textbf{REJECTED}. \texttt{coverageMode}
remains semantically required.

The mode does not require one serialization; an unambiguous textual
equivalent is sufficient.

\subsection{4. Attack B - retain omissionSemantics as an independent
field}\label{attack-b---retain-omissionsemantics-as-an-independent-field}

T2 carried both \texttt{coverageMode} and \texttt{omissionSemantics}.

T3 asks whether the latter adds independent necessary semantics.

Once mode semantics are fixed:

\begin{Shaded}
\begin{Highlighting}[]
\NormalTok{EXHAUSTIVE\_FOR\_DECLARED\_SCOPE}
\NormalTok{  omitted eligible+qualified item {-}\textgreater{} projection defect}

\NormalTok{SELECTIVE}
\NormalTok{  omitted eligible item {-}\textgreater{} no project{-}semantic conclusion}
\end{Highlighting}
\end{Shaded}

an independent omission policy can only duplicate those rules or
contradict them.

\textbf{Disposition:} \texttt{omissionSemantics} as an independent
mandatory field \textbf{REJECTED AS REDUNDANT}.

This is the only material simplification accepted by the closure review.

\subsection{5. Attack C - collapse qualificationPolicy into generic
coverage
text}\label{attack-c---collapse-qualificationpolicy-into-generic-coverage-text}

A scope such as \texttt{transfer\ propositions} does not state whether
stale, diagnostic or pending material is admissible.

A current-state projection and a review projection can have the same
semantic topic but intentionally different admissible BA state.

If the distinction is hidden only in rendering conventions, a non-visual
consumer can ingest stale/unresolved basis as if current.

\textbf{Disposition:} semantic collapse \textbf{REJECTED}.

\texttt{qualificationPolicy} may be physically nested in coverage
metadata, but its responsibility remains distinct.

\subsection{6. Attack D - remove
interpretationRuleRef}\label{attack-d---remove-interpretationruleref}

The T2 overlapping-rule counterexample is replayed:

\begin{Shaded}
\begin{Highlighting}[]
\NormalTok{R1: serviceUse + responsibilityBoundary}
\NormalTok{    {-}\textgreater{} external{-}service exposure}

\NormalTok{R2: serviceUse + responsibilityBoundary + constraints}
\NormalTok{    {-}\textgreater{} assurance dependency review}
\end{Highlighting}
\end{Shaded}

An item with only a BA trace does not identify which local rule produced
its meaning when several rules overlap.

\textbf{Disposition:} removal \textbf{REJECTED}.

\texttt{interpretationRuleRef} remains required for meaning-bearing
method-owned interpretation. It remains projection-local and never
becomes a BA3 derivation rule.

\subsection{7. Attack E - make trace roles mandatory
everywhere}\label{attack-e---make-trace-roles-mandatory-everywhere}

A single BA element mapped directly to one human bullet does not need an
artificial contribution role merely to satisfy serialization uniformity.

However, role distinctions are essential when multiple BA inputs
contribute differently to one interpretation.

\textbf{Disposition:} universal mandatory role key \textbf{REJECTED};
conditional role binding \textbf{ACCEPTED}.

Final rule:

\begin{Shaded}
\begin{Highlighting}[]
\NormalTok{traceRoleKey required where contribution roles differ}
\end{Highlighting}
\end{Shaded}

\subsection{8. Attack F - collapse shared rendering and method-owned
interpretation}\label{attack-f---collapse-shared-rendering-and-method-owned-interpretation}

\subsubsection{F1 semantic-preserving
aggregate}\label{f1-semantic-preserving-aggregate}

\begin{Shaded}
\begin{Highlighting}[]
\NormalTok{C constraint}
\NormalTok{I constraint}
\NormalTok{provenance constraint}
\NormalTok{    {-}\textgreater{} enumerated "delivery constraints" group}
\end{Highlighting}
\end{Shaded}

The group is no stronger than its BA basis and can remain shared
rendering.

\subsubsection{F2 consumer conclusion}\label{f2-consumer-conclusion}

\begin{Shaded}
\begin{Highlighting}[]
\NormalTok{same constraint basis}
\NormalTok{    {-}\textgreater{} "high{-}assurance channel"}
\end{Highlighting}
\end{Shaded}

This is a stronger consumer interpretation. If admissible at all, it
requires a method-owned rule and cannot be presented as shared BA
meaning.

\textbf{Disposition:} collapse \textbf{REJECTED}.

The semantic distinction remains even if one physical item contains both
a shared rendering and local annotations.

\subsection{9. Attack G - force one shared projection
ontology}\label{attack-g---force-one-shared-projection-ontology}

T3 replays three projections over the same BA:

\begin{enumerate}
\def\labelenumi{\arabic{enumi}.}
\tightlist
\item
  a human project-understanding view;
\item
  a flow/exposure method view; and
\item
  an assurance/responsibility method view.
\end{enumerate}

The method views use intentionally incompatible local categories. Both
can remain valid because their common denominator is role-bound BA
trace, not category equivalence.

The human view can group the same BA basis without adopting either
method taxonomy.

\textbf{Disposition:} universal projection ontology \textbf{REJECTED}.

No \texttt{SharedProjectionConcept} or method-category equivalence
relation is forced.

\subsection{10. Attack H - require cross-baseline projection-item
identity}\label{attack-h---require-cross-baseline-projection-item-identity}

M1-M4 are replayed under rebuild semantics.

A local projection item at B0 may disappear, split or be represented
differently at B1 because the BA meaning or the projection descriptor
changed.

The project-semantic continuity already exists in BA3. Adding
\texttt{RETAIN\ \textbar{}\ REPLACE\ \textbar{}\ RETIRE} to projection
items duplicates authority without current evidence.

\textbf{Disposition:} projection-item BA-like lifecycle \textbf{REJECTED
AS UNNECESSARY}.

Method/UI continuity remains tool-owned if desired.

\subsection{11. Integrated rebuild
regression}\label{integrated-rebuild-regression}

\subsubsection{M1 Ethernet -\textgreater{}
Wi-Fi}\label{m1-ethernet---wi-fi}

Flow projection may remain materially stable while realization/assurance
projection changes. PASS: different projection deltas are expected
because coverage differs.

\subsubsection{M2 external -\textgreater{} project-owned
transport}\label{m2-external---project-owned-transport}

The external-service exposure interpretation becomes non-applicable when
the BA responsibility/service basis changes. A different method can
produce a project-owned authority item. PASS: no common method category
is required.

\subsubsection{M3 remote -\textgreater{} local
recognition}\label{m3-remote---local-recognition}

Retired remote transfer disappears from a current exhaustive flow view.
Unrelated responsibility items can remain. PASS: rebuild follows BA
rather than local item history.

\subsubsection{M4 raw capture -\textgreater{} opaque reference
conflict}\label{m4-raw-capture---opaque-reference-conflict}

A current projection excludes stale/unresolved transfer meaning; a
review projection may surface it with qualification. PASS: no projection
invents post-resolution project truth.

\subsection{12. Order/WMS regression}\label{orderwms-regression}

Externalizing inventory authority can produce large deltas in an
integration projection and different deltas in a
responsibility/assurance projection.

Each remains consistent if it follows its own declared scope, mode and
qualification policy and traces to accepted BA.

No method-local \texttt{ExternalSystem}, \texttt{TrustBoundary} or
\texttt{IntegrationRisk} category is forced into BA.

\textbf{Disposition:} PASS.

\subsection{13. Provider-normalization
regression}\label{provider-normalization-regression}

A consumer needing \texttt{PaymentAuthorizationResult} must consume
accepted normalized BA meaning.

The negative pattern remains rejected:

\begin{Shaded}
\begin{Highlighting}[]
\NormalTok{provider raw vocabulary}
\NormalTok{    {-}\textgreater{} hidden method mapping}
\NormalTok{        {-}\textgreater{} supposedly shared project result}
\end{Highlighting}
\end{Shaded}

The accepted path remains:

\begin{Shaded}
\begin{Highlighting}[]
\NormalTok{provider raw vocabulary}
\NormalTok{    {-}\textgreater{} governed mapping / BA3 derivation}
\NormalTok{        {-}\textgreater{} accepted BA meaning}
\NormalTok{            {-}\textgreater{} projection interpretation}
\end{Highlighting}
\end{Shaded}

\textbf{Disposition:} PASS.

\subsection{14. Orthogonality check}\label{orthogonality-check}

The closure review confirms four independent questions:

\begin{Shaded}
\begin{Highlighting}[]
\NormalTok{trace}
\NormalTok{  Which BA meaning supports this projection item?}

\NormalTok{coverage}
\NormalTok{  Is the view selective or exhaustive for its declared eligible scope?}

\NormalTok{qualification}
\NormalTok{  Which BA states are admissible, and how is non{-}current meaning preserved?}

\NormalTok{interpretation rule}
\NormalTok{  Which consumer{-}owned rule justifies this non{-}trivial local meaning?}
\end{Highlighting}
\end{Shaded}

No pair can be merged without losing a reviewed distinction.

Physical co-location remains allowed.

\subsection{15. Earlier-layer reopen
checks}\label{earlier-layer-reopen-checks}

No material counterexample forces an earlier reopen.

\begin{Shaded}
\begin{Highlighting}[]
\NormalTok{BA1: BAReferent + BAProposition remain sufficient       PASS}
\NormalTok{BA2: current operator/role vocabulary remains sufficient PASS}
\NormalTok{BA3: provenance/lifecycle/change continuity sufficient  PASS}
\end{Highlighting}
\end{Shaded}

A method-specific node, edge, risk class, assurance class, UI state or
local identifier is not a shared-semantic counterexample.

\subsection{16. Closure decision}\label{closure-decision}

\begin{Shaded}
\begin{Highlighting}[]
\NormalTok{BA4{-}T3   CLOSED / PASS}
\NormalTok{BA4      CLOSED FOR CURRENT THESIS SCOPE}
\end{Highlighting}
\end{Shaded}

The accepted final contract is:

\begin{itemize}
\tightlist
\item
  \texttt{BA4\_PROJECTION\_BOUNDARY\_TRACEABILITY\_INTERPRETATION\_COVERAGE\_CONTRACT\_R1.md}
\end{itemize}

BA4-T1/T2 artifacts remain immutable research derivation history.

\subsection{17. Closure ledger summary}\label{closure-ledger-summary}

\begin{Shaded}
\begin{Highlighting}[]
\NormalTok{projection authority                                    REJECTED}
\NormalTok{immutable projection revision                           RETAIN}
\NormalTok{eligibleBAScope                                         RETAIN}
\NormalTok{coverageMode                                            RETAIN}
\NormalTok{independent omissionSemantics                           REMOVE / REDUNDANT}
\NormalTok{qualificationPolicy                                     RETAIN}
\NormalTok{BA trace                                                RETAIN}
\NormalTok{traceRoleKey mandatory everywhere                       REJECT}
\NormalTok{traceRoleKey where roles differ                         RETAIN}
\NormalTok{shared semantic rendering                               RETAIN}
\NormalTok{method{-}owned interpretation                             RETAIN}
\NormalTok{interpretationRuleRef                                   RETAIN WHEN APPLICABLE}
\NormalTok{universal projection DSL                                REJECT}
\NormalTok{universal projection ontology                           REJECT}
\NormalTok{cross{-}projection shared identity                        NOT FORCED}
\NormalTok{BA trace + BA3 continuity common anchor                 ACCEPT}
\NormalTok{projection lifecycle ontology                           REJECT}
\NormalTok{new BAE family                                          NOT FORCED}
\NormalTok{new BA2 operator                                        NOT FORCED}
\NormalTok{BA1/BA2/BA3 reopen                                      NOT TRIGGERED}
\end{Highlighting}
\end{Shaded}

\subsection{18. BA4 reopen criteria}\label{ba4-reopen-criteria}

Reopen only the smallest responsibility on a material case where the
final contract cannot express necessary projection completeness, cannot
preserve qualification, cannot reproduce method-owned interpretation,
loses source drill-down, requires a shared semantic identity beyond BA
trace, or cannot rebuild from BA3 continuity without a second shared
lifecycle.

Do not reopen for ordinary method taxonomy growth, UI identity needs,
renderer changes, new projection revisions or storage preferences.

\subsection{19. Next phase boundary}\label{next-phase-boundary}

BA5 may begin only after this BA4 closure package is committed, pushed
and remotely verified.

BA5 must not reinterpret projection-local labels as BA semantic keys.
Its first trial must pressure the boundary among source wording, stable
BA semantic keys, projection display labels and authoring synonyms
before optional assistance is considered.

\end{document}
